\usepackage[utf8]{inputenc}
\usepackage[T1]{fontenc}

\usepackage{amsmath,amssymb}
\usepackage{graphicx}
\usepackage{booktabs}
\usepackage{longtable}
\usepackage{listings}
\usepackage{xcolor}
\usepackage[hidelinks]{hyperref}
\lstset{basicstyle=\ttfamily\small,columns=fullflexible,keepspaces=true,
  frame=single,breaklines=true}
% names
\newcommand{\ac}{AnharmoniCaOs}
\newcommand{\acversion}{2026.11}
% an input option, a file name, a program name
\newcommand{\opt}[1]{\texttt{#1}}
\newcommand{\file}[1]{\texttt{#1}}
\newcommand{\prog}[1]{\texttt{#1}}
% pointer to the theory notes ("Anharmonic oscillations with the Van Vleck
% method"); #1 is the section or equation, as text
\newcommand{\notes}[1]{(theory notes, #1)}
% marks what exists only in the development version (dev/polyad-jobs),
% not in the release v2026.11
\newcommand{\devonly}{\textsf{[dev]}}
% marks a feature that does not work reliably yet and is not recommended
\newcommand{\wip}{\textsf{[work in progress]}}
% an option-reference entry: name, default, phase, description
\newcommand{\optentry}[4]{\item[\opt{#1}] \emph{default} \texttt{#2}, \emph{phase} #3.\ #4}
